\title{LongRoPE: Extending LLM Context Window Beyond 2 Million Tokens}

\begin{abstract}
Expanding the context window in large language models (LLMs) is highly sought after. Nevertheless, the high costs of fine-tuning, the limited availability of lengthy textual data, and the emergence of problematic values from unfamiliar token positions restrict current large context windows to approximately 128k tokens.
\end{abstract}

In this work, we present LongRoPE, a novel approach that, for the first time, extends the context window of pre-trained LLMs to a remarkable \textbf{2048k} tokens. This is accomplished with as few as 1k fine-tuning steps at training lengths up to 256k, while preserving the original model's performance on shorter context windows. Our methodology is centered on three main innovations: (i) we detect and leverage two distinct types of non-uniformities in positional interpolation through an optimized search process, yielding improved fine-tuning initialization and enabling up to an 8$\times$ extension even without additional fine-tuning; (ii) we propose a staged extension procedure, whereby the LLM is initially fine-tuned to a 256k context, followed by a second positional interpolation on this already-extended model, ultimately resulting in a context window of 2048k; (iii) we recalibrate LongRoPE on sequences of length 8k, effectively restoring the model’s capability in shorter context settings. Comprehensive evaluations using LLaMA2 and Mistral over a wide range of tasks validate the efficacy of our approach. Models adapted via LongRoPE maintain their original architecture, with only minimal changes to the positional embedding layer, and preserve compatibility with existing optimizations. The codebase will be made accessible at \texttt{https://github.com/microsoft/LongRoPE}
\end{abstract}

\section{Introduction}

Large Language Models (LLMs) have demonstrated impressive capabilities across a range of applications~\cite{gpt4,llama2}, yet they are constrained by a finite context window size—LLaMA2, for instance, is limited to 4096 tokens~\cite{llama2}. When input exceeds this context window, the quality of the model's outputs deteriorates because the model lacks training for handling tokens at these extended positions. Such limitations create obstacles in critical use cases, such as in-context learning involving many examples~\cite{Huang2023FewerIM}, as well as for LLM-based agents~\cite{park2023generative,madaan2023selfrefine}.

Recent studies have demonstrated that the context window of a pre-trained LLM can be expanded to approximately 128k through fine-tuning with extended text sequences~\cite{longlora,pi,yarn,zhang2024soaring,liu2023scaling}. Nonetheless, further increasing the context window faces three significant challenges. Firstly, the inclusion of previously unseen position indices, which remain untrained, introduces anomalous values; this results in out-of-distribution complications and hinders convergence during fine-tuning~\cite{pi}. This issue is especially acute when scaling the context window from 4k to values above $>$1000k, where over 90\% of the position indices are novel. Secondly, effective fine-tuning relies on access to textual data that matches the intended longer lengths. However, data sources with extensive sequences—especially those surpassing 1000k tokens—are scarce in existing datasets. Additionally, training on such lengthy samples incurs substantial computational costs, demanding excessive time and GPU resources. Thirdly, at these extreme context lengths, the model’s attention distribution becomes excessively diffused across the many possible token positions, which in turn undermines the model’s accuracy on tasks with shorter input lengths~\cite{pi}.

A common strategy to address the initial issue involves interpolating RoPE positional embeddings~\cite{rope,pi}, a method that scales novel positional indices back into the domain on which the model was initially trained, as illustrated in Fig.\ref{fig:overview}. The Position Interpolation (PI)~\cite{pi} technique achieves this by applying a linear interpolation to the rotary angles of RoPE, guided by the extension ratio. Alternatively, NTK~\cite{ntk,dynamicntk} proposes adopting distinct strategies of interpolation and extrapolation depending on the RoPE dimension. YaRN~\cite{yarn} further organizes the RoPE dimensions into three frequency groups and selectively applies extrapolation, NTK, or linear interpolation to each group. Despite these interventions, the information entropy encoded by positional embeddings in Transformer architectures is \emph{highly non-uniform and intricate}. Current methodologies fail to sufficiently exploit this nuanced non-uniformity, resulting in a loss of information that constrains the achievable context window length.

Section~\ref{sec:analysis} presents two principal empirical observations: \textbf{(1)} Successful positional interpolation must account for two specific sources of non-uniformity—differences across RoPE dimensions and across token positions. Reduced interpolation proves advantageous for smaller RoPE dimensions and for tokens closer to the sequence beginning; nevertheless, the most suitable degree of interpolation varies with the intended length of the extension.
\textbf{(2)} Incorporating these non-uniformities into the interpolation scheme makes it possible to better preserve the critical information encoded in the original RoPE, especially within important dimensions and at key token positions. This approach significantly limits information loss stemming from interpolation and thereby creates stronger initializations for fine-tuning. Additionally, it enables up to an 8$\times$ sequence extension, even in the absence of fine-tuning.

Building upon these results, we propose \textbf{LongRoPE}, a robust approach designed to expand the LLM context window to over 2 \emph{million} tokens. The architecture of LongRoPE integrates three principal advancements. The first centers on fully leveraging the multidimensional non-uniform characteristics inherent in positional interpolation. Specifically, LongRoPE discovers optimal rescaling factors for the rotation angles associated with RoPE in each dimension, guided by the locations of the tokens. Given that the search space for identifying these rescaling factors grows exponentially with the desired extension ratio, LongRoPE employs an evolutionary search strategy, augmented by two specialized optimization methods, to enhance the search's effectiveness. An illustrative instance of the resulting rescaled RoPE is depicted in Fig.~\ref{fig:overview}.

Next, {LongRoPE} adopts an efficient, gradually expanding approach to reach a 2048k context window without necessitating direct fine-tuning on extremely lengthy texts, which are limited and difficult to obtain. The process initiates with a search for a suitable 256k context length on the pre-trained LLM, followed by fine-tuning at this identified length. Subsequently, leveraging our non-uniform positional interpolation—which supports up to an 8$\times$ extension without further fine-tuning—we carry out another search to determine new RoPE rescale factors on the previously fine-tuned and expanded LLM. Through this two-phase process, a 2048k context window is ultimately realized for both LLaMA2 and Mistral~\cite{mistral}.

In order to prevent a drop in performance on the initial (shorter) context windows, {LongRoPE} maintains adjustment of the RoPE rescale factors within the expanded LLM. Analogous to the process of increasing context length from 256k to 2048k, the model is also downscaled to 4k and 8k context windows on the 256k fine-tuned LLM, utilizing our search algorithm to minimize positional interpolation. When performing inference, if the input sequence falls below 8k tokens, RoPE is modified with the rescale factors identified through our search process.

A comprehensive set of experiments involving multiple LLM architectures and a range of long-context tasks provides evidence for the efficacy of our approach. Our results indicate that LongRoPE consistently preserves low perplexity across evaluation lengths from 4k up to 2048k, secures more than 90\% accuracy in passkey retrieval tasks, and attains performance on par with existing models on standard benchmarks built within a 4096-token context window. Furthermore, LongRoPE is compatible with any LLM employing RoPE embeddings. We plan to make our implementation and LongRoPE-2048k models publicly available.

\section{Non-uniformity in Positional Interpolation}
\label{sec:analysis}

\subsection{Preliminary}

Transformer models require explicit positional information, often in the form of position embedding, to represent the order of input tokens. Our work focuses on the RoPE~\cite{rope} position embedding, which is widely used in recent  LLMs. For a token at position index $n$, its corresponding RoPE encoding can be simplified as follows:

\begin{equation}
		\fontsize{7.8}{7.8} \selectfont
	\label{eq:rope}
	\left[ cos(n\theta_0), sin(n\theta_0), cos(n\theta_1), \cdot\cdot\cdot, cos(n\theta_{d/2-1}), sin(n\theta_{d/2-1}) \right]   
\end{equation}

where $d$ denotes the embedding dimension, 
$n\theta_i$ specifies the rotary angle for a token located at position $n$, and 
$\theta_i=\theta^{-2i/d}$ determines the frequency of rotations. RoPE sets the default base value of $\theta$ to 10000.

\textbf{\emph{Context window extension ratio $s$} and positional interpolation}. We characterize $s$ as the proportion between the new, extended context window length $L'$ and the initial length $L$, given by $s = \frac{L'}{L}$.

To extend the context window from $L$ to $L'$, current positional interpolation methods suggest downscaling rotation frequencies $\theta_i$ based on extension ratio $s$. Let $\beta=\theta^{2/d}$, and $\lambda$ denote the actual rescale factor related to $s$, we   unify these positional interpolation methods as follows:

\begin{equation}	
	\fontsize{7.5}{7.5} \selectfont
	\label{eq:unified_rope}
	\begin{aligned}
		\left[cos\left(\frac{n}{\lambda(\beta)^0}\right), sin \left(\frac{n}{\lambda(\beta)^0}\right), cos\left(\frac{n}{\lambda(\beta)^1}\right), \cdot\cdot\cdot,  sin \left(\frac{n}{\lambda(\beta)^{d/2-1}}\right) \right]   
	\end{aligned}
\end{equation}

\textbf{Linear positional interpolation (PI)}. PI~\cite{pi} implements a straightforward linear rescaling of positional indices, constrained within the pre-trained sequence length. When extending to a target ratio $s$, each rotation angle across RoPE dimensions is scaled down uniformly by $\lambda=s$. This approach compresses positional encoding, causing the position representation to become congested and impairing the model's discrimination between nearby tokens. As a consequence, PI exhibits diminished performance as the extension ratio increases.

\textbf{NTK-based interpolation and extrapolation}. Studies such as~\cite{ntk,dynamicntk} reframe RoPE through the lens of information encoding and employ Neural Tangent Kernel (NTK) theory~\cite{jacot2018neural,tancik2020fourier} for analysis. To address the positional crowding problem found in PI, they advocate for redistributing interpolation "pressure" among the RoPE dimensions: specifically, scaling down the influence on lower (high-frequency) dimensions while scaling up on higher (low-frequency) ones. This mechanism enables both positional interpolation and extrapolation, parameterized by $\lambda=s^{i}$. The dynamic NTK variant~\cite{dynamicntk} further enhances this approach by modifying the extension ratio per position in accordance with the sequence length. In contrast to PI, which demands model fine-tuning, NTK-aware approaches support context window extension without such fine-tuning, although the achievable extension ratio is typically limited to 4$\times$.

\textbf{YaRN}~\cite{yarn} proposes a notable advancement in positional interpolation by segmenting RoPE dimensions into three distinct categories according to their associated frequencies, and then applying separate interpolation methods to each group. Specifically, dimensions with high frequencies are extrapolated ($\lambda$=1), whereas those with low frequencies are interpolated linearly (PI). The remaining, intermediate-frequency dimensions are assigned the NTK method. The effectiveness of YaRN is primarily attributed to this frequency-based classification of RoPE dimensions; however, the determination of these groups currently relies on manual, empirical approaches, which can potentially limit performance when adapting to novel LLMs.

\subsection{Study on Non-uniform Positional Interpolation}

Drawing upon insights from NTK and YaRN, we observe that their improvements stem from introducing non-linearity—namely, through treating distinct RoPE dimensions with frequency-dependent mechanisms for tailored interpolation and extrapolation. Yet, the prevailing non-linear approaches are predominantly based on heuristics crafted by humans. This leads to two central questions: (1) Does the current approach to positional interpolation represent the optimal solution? (2) Might there exist non-linearities that have yet to be investigated?

In order to address these inquiries, we employ evolution search (refer to Sec.~\ref{sec:method}) to identify improved non-uniform positional interpolation strategies for LLaMA2-7B. This process is steered by perplexity measurements, evaluated across 5 randomly selected samples from the PG19~\cite{pg19} validation set. Our empirical investigation yields the following principal insights.

\textbf{Finding 1}: \textit{RoPE dimensions exhibit substantial non-uniformities, which are not effectively handled by current positional interpolation methods.}

We identify the optimal value of $\lambda$ for each RoPE dimension as specified in Eq.~\ref{eq:unified_rope}. In Table~\ref{tbl:exp1}, we present a perplexity comparison of various methods, including ours, for the LLaMA2-7B model evaluated on the PG19 and Proof-pile~\cite{proof-pile} test datasets, all without any additional fine-tuning. The approach derived from our search achieves notable improvements in performance, indicating that both the linear (PI) and non-uniform (Dynamic-NTK and YaRN) interpolation strategies do not represent optimal solutions. It is also worth noting that YaRN lags behind both PI and NTK on PG19, as it fails to attain the intended context window length when applied to language models that have not been fine-tuned. Specifically, the perplexity for YaRN rises sharply beyond 7k within an 8k context window.

Our exploration reveals that the rescaling coefficients $\lambda$ in Eq.~\ref{eq:unified_rope} are non-uniform, in contrast to the constant scaling $s$ used in PI, the analytical computation in NTK, and the group-based scheme in YaRN. The adoption of these non-uniform scaling factors leads to a notable enhancement in LLaMA2's language modeling abilities—as indicated by reduced perplexity—across 8k and 16k context lengths, achieved without any additional fine-tuning. The primary reason for this improvement is that the resultant positional embeddings retain the essential properties of the original RoPE, particularly for crucial dimensions, thereby alleviating the challenge for large language models in accurately differentiating nearby token positions.

\textbf{Finding 2}: \textit{RoPE for the initial tokens in the input sequence should be extrapolated with less interpolation.} 

We propose that for the first $\hat{n}$ tokens of input sequences, position interpolation applied by RoPE should be minimized. The rationale is that these initial tokens typically attain high attention weights, making them particularly significant for the attention mechanism, as demonstrated in works like Streaming LLM~\cite{streamingllm} and LM-Infinite~\cite{han2023lminfinite}. To examine this hypothesis, we extend context windows to sizes of 8k and 16k using PI and NTK methods, while systematically withholding interpolation for the first $\hat n$ (where $\hat n$ = 0, 2, ..., 256) tokens. Setting $\hat n = 0$ reverts to the default PI and NTK configurations. The results, summarized in Table~\ref{tbl:tokenposition}, reveal two main findings: \textbf{(1)} exempting the initial tokens from position interpolation positively impacts performance; \textbf{(2)} the ideal value of $\hat n$ varies according to the desired context window length.

\textbf{Finding 3}: \textit{Non-uniform positional interpolation effectively extends LLM context window in both fine-tuning and non-fine-tuning settings.} 

Our experiments demonstrate that implementing our optimized non-uniform position interpolation yields marked improvements in extension capabilities at the 8k and 16k context lengths without the need for fine-tuning. However, for even longer context lengths, such as those beyond 16k, fine-tuning becomes necessary. Therefore, we conduct fine-tuning on LLaMA2-7B using our optimized RoPE scheme for an extended 64k context window (refer to the Appendix for detailed configuration). As evidenced in Table~\ref{tbl:finetune}, this approach delivers substantially better results than both PI and YaRN, outperforming these methods before and after fine-tuning the LLaMA2-7B model. These enhancements arise from our strategic application of non-uniform positional interpolation, which limits information degradation and establishes a superior starting point for subsequent fine-tuning.

\textbf{Summary}. Our analysis identifies two key forms of non-uniformity: disparities across RoPE dimensions and across token positions. Effectively harnessing these non-uniform characteristics in positional interpolation techniques leads to significant gains in LLM context extension capabilities.

\section{LongRoPE}

\label{sec:method}
Building on these insights, we propose LongRoPE. This approach initially deploys a resourceful search algorithm to comprehensively leverage both forms of non-uniformity, subsequently utilizing this mechanism to expand the LLM context window to exceed 2 million tokens.

\subsection{Problem Formulation}

These two forms of non-uniformity expand the range of possible solutions and add layers of complexity to the optimization process. To tackle this issue, we conceptualize the multidimensional non-uniform position interpolation optimization as a search-based problem.

For a LLM targeting a  context window size of $L'$ and lengthy input documents $\mathbf{X}$, where each $\mathbf{x}\in\mathbf{X}$ surpasses $L'$ in token length, we denote the original rotary angle of the $i^{th}$ dimension in RoPE embedding at token position $n$ as  $\frac{n}{\beta^i}$. The optimization problem is then formulated as follows:

\begin{equation}
\begin{gathered}
		\label{eq:problem}

\underset {\mathbf{x}\in\mathbf{X}; \, |\mathbf{x}|\ge L'}{ \operatorname {arg\,min} } \,  \mathcal{L}\, \left( \text{LLM} (\text{RoPE}, \mathbf{X})\right), \text{where}
\\
\scalebox{0.93}{
$\underset {\begin{subarray}{l} i=0,\cdot\cdot,\frac{d}{2}-1;\\ n\in[ 0,|\mathbf{x}|);\end{subarray}}{\text{RoPE}_(n)}= {\left[\cdot\cdot,\cos\left(\mathbb{I}(\hat{\lambda}_{i}, \hat{n})\cdot\frac{n}{\beta^i}\right), \sin\left(\mathbb{I}(\hat{\lambda}_{i}, \hat{n})\cdot\frac{n}{\beta^i}\right),\cdot\cdot\right]}$}
\\
\text{where} \quad \mathbb{I}(\hat{\lambda}_{i},\hat{n}) = 
\begin{cases}
1 & \text{if } n < \hat{n} \\
\frac{1}{\lambda_{i}} & \text{if } n \geq \hat{n}
\end{cases}
\end{gathered}
\end{equation}

In this approach, we define a set of scaling factors $\mathbb{I}(\hat{\lambda}_{i},\hat{n})$ to address two types of non-uniformities. Here, $\hat{\lambda}_i$ captures irregularities across RoPE dimensions, while $\hat{n}$ represents position-based non-uniformity among tokens. The function $\mathbb{I}(\hat{\lambda}_{i},\hat{n})$ modifies the rotation angle for the $i^{th}$ RoPE dimension, using $\hat{\lambda}_i$ as the scaling parameter and $\hat{n}$ as the positional cutoff. For token positions less than $\hat{n}$, i.e., up to $\hat{n}-1$, the scaling parameter $\hat{\lambda}_i$ does not modify the computation, and the standard RoPE rotary angle $\frac{n}{\beta^i}$ remains in use. When the token position satisfies $n \geq \hat{n}$, the scaling factor is employed.

Suppose a desired context window of size $L'$ is specified; our aim is to determine the optimal set of rescale factors ($\mathbb{I}(\hat{\lambda}_{0},\hat{n})$, $\mathbb{I}(\hat{\lambda}_{1},\hat{n})$, ..., $\mathbb{I}(\hat{\lambda}_{i},\hat{n})$, ...) spanning from the $1^{st}$ up to the $d^{th}$ RoPE dimension. By applying these rescaling adjustments to the target $\text{LLM}$'s $\text{RoPE}$ mechanism, we seek to minimize the next-token prediction loss, $\mathcal{L}$ (corresponding to perplexity), across input sequences $\mathbf{X}$ of length $L'$.

\subsection{Searching the Non-uniform Position Interpolation}

To address the issue outlined in Eq.~\ref{eq:problem}, we propose a straightforward but powerful approach that seeks the best RoPE rescale factors, thereby making full use of the multidimensional variations present in the position embedding.

\textbf{Search space}. We construct an extensive search space that accommodates the two types of non-uniformities. The structure of this search space is presented in Table~\ref{tbl:searchspace}. In particular, our approach enables independent optimization of a rescale factor for each RoPE embedding dimension. To streamline the search process, we opt to search over $\lambda_i$ and $\hat{n}$, rather than directly searching the composite function $\mathbb{I}(\hat{\lambda}_{i},\hat{n})$, with the relation $\hat{\lambda}_{i}=1/\lambda_i$. As detailed in Table~\ref{tbl:searchspace}, the search for $\lambda_i$ spans from a lower bound of 1.0 (which equates to straightforward extrapolation) up to an upper bound of $s\times1.25$ (representing an interpolation wider than PI), using increments of 0.01, where $s$ denotes the desired context window scaling factor.

The parameter $\hat{n}$ determines how many of the earliest token positions utilize the standard RoPE embedding rather than undergoing position interpolation. In our experiments, we vary $\hat{n}$ over the set \{0, 1, 2, 4, 8, 12, 16, 20, 24, 28, 32, 64, 128, 256\}. Setting $\hat{n}=0$ indicates that every token position is assigned the computed rescale factors.

\textbf{Evolution-based search}. The search space outlined in Table~\ref{tbl:searchspace} encompasses a vast array of positional interpolation options, rendering efficient navigation highly challenging. For instance, applying a $s=4\times$ extension results in $400^{128/2}\times14 = 4\times10^{167}$ possible configurations. As the extension ratio increases, the size of the search space grows exponentially. To manage this complexity, we employ evolution search~\cite{guo2020single} and incorporate two optimization strategies to significantly enhance the efficiency of the search process. The general framework of the search algorithm is provided in Algorithm~\ref{alg:search}.

\textit{Optimized initial population generation}. In lieu of purely random initialization for a population of $P$ rescale factors, our approach includes explicitly adding the three RoPE rescale factors associated with PI, NTK, and YaRN as members of the initial population. For the other $P-3$ individuals, we introduce diversity by applying random mutations to these three baseline rescale factors, each with a probability $p$.

\textit{Monotonically non-decreasing constraint}. Once the initial population has been created, we evaluate the LLM perplexity associated with each candidate. This involves applying each individual’s specific RoPE rescale factors to the target LLM and subsequently determining the perplexity for input $\mathbf{X}$. The best $k$ candidates are then selected as parents for the subsequent evolutionary steps. Due to the vastness of the search space, simple mutation and crossover operations can frequently result in suboptimal candidates, leading to many redundant perplexity evaluations. This inefficiency is exacerbated when $L'$ is large, as each perplexity inference demands considerable computational effort.

To overcome this issue, we enforce a monotonic non-decreasing condition on the selected RoPE rescaling factors, formalized as $\lambda_i \le \lambda_{i+1}$. Only those RoPE configurations adhering to this criterion are utilized for perplexity assessment on LLMs, which leads to a notable reduction in search complexity. In particular, we stipulate that the sequence $\lambda_i$ must grow monotonically with respect to the RoPE dimension ($i=0,...,63$). This requirement of monotonicity across dimensions is motivated by the insights from NTK theory~\cite{jacot2018neural,tancik2020fourier,ntk}: it suggests that dimensions associated with higher frequency (i.e., lower $i$) necessitate minimal interpolation, thus favoring smaller values of $\lambda_i$, whereas dimensions related to lower frequency (i.e., higher $i$) can tolerate greater interpolation, corresponding to larger $\lambda_i$ values.

\begin{algorithm}[t]
	\small
	\caption{The search algorithm}
	\textbf{Input:} target LLM, input samples $\mathbf{X}$, population size $P$, mutation size $N_1$, crossover size $N_2$, maximum iterations $\mathcal{T}$, mutation probability $p$\\

	\begin{algorithmic}[1]
		\label{alg:search}
		\STATE Initialize Top-k as the empty set;	
		\STATE Set $\text{P}_0$ by calling \textit{Initialize\_population\_with\_optimization} ($P$, $\mathbf{X}$, $p$); 
		\FOR{$i$ from $1$ to $\mathcal{T}$}
		\STATE Invoke \textit{Compute\_perplexity} (LLM, $\text{P}_{i-1}$, $\mathbf{X}$);
		\STATE  Update Top-k with \textit{Update\_Topk} (Top-k, $\text{P}_{i-1}$);
		\STATE Generate $\text{P}_{mutation}$ using \textit{Mutation\_with\_mono\_constraint} (Top-k, $N_1$, $p$); 
		\STATE Produce $\text{P}_{crossover}$ by calling \textit{Crossover\_with\_mono\_constraint} (Top-k, $N_2$);
		\STATE Form the next population $\text{P}_i$ as the union of $\text{P}_{mutation}$, $\text{P}_{crossover}$, and Top-k;
		\ENDFOR
		\STATE Output the member of Top-k achieving minimal perplexity;
	\end{algorithmic}
\end{algorithm}

\textbf{8$\times$ extension without fine-tuning}. Utilizing evolutionary search, we efficiently discover optimal non-uniform RoPE rescale factors, which maintain essential positional and dimensional components to reduce information degradation due to interpolation. As illustrated in Fig.\ref{fig:non-ft}, this approach successfully enlarges LLaMA2’s context window from 4k to 32k tokens without the need for fine-tuning. By comparison, alternative approaches like PI, and non-uniform NTK and YaRN experience a significant increase in perplexity when the context window is doubled.

\subsection{Extending LLM Context Window to 2048K}
\label{sec:extendto2048k}

\textbf{Progressive extension to 2048k}. In this section, we present our approach for expanding the context window of existing pre-trained LLMs from the standard 4k up to more than 2048k. As previously shown, our non-uniform positional interpolation enables up to 8$\times$ context extension without necessitating any fine-tuning. However, for much larger increases (such as a 512$\times$ extension), fine-tuning becomes essential. One strategy involves identifying appropriate RoPE rescaling factors tailored for the target 2048k window and subsequently applying fine-tuning. Nonetheless, this approach encounters significant obstacles, as the resource demands for such extensive training are exceedingly high. Additionally, we have observed in practice that effectively fine-tuning LLMs with such a large extension ratio proves to be quite difficult (refer to Appendix for details).

Luckily, LongRoPE functions well for both the base model and the further extended LLMs after fine-tuning. As a result, we propose an incremental and efficient approach that reaches the target of 2048k using only 1k fine-tuning iterations, all conducted within a training context length of 256k.

$\Diamond$ \textit{LongRoPE search for expanding pre-trained LLMs to 256k.} Using LLaMA2 for illustration, we perform a search targeting context window sizes of 128k and 256k. At this phase, the extension ratios achieved are 32$\times$ and 64$\times$, respectively.

$\Diamond$ \textit{Fine-tuning to 256k}. Subsequently, we adapt the pre-trained LLM to support a 256k context window through staged fine-tuning. Initially, we perform 400 fine-tuning steps on LLaMA2 utilizing RoPE rescaled factors designed for 128k. After this phase, we substitute these factors with ones corresponding to 256k on the resultant checkpoint and continue with an additional 600 fine-tuning steps. This incremental approach demonstrates superior efficiency compared to attempting a direct fine-tuning to 256k.

$\Diamond$ \textit{Scaling up fine-tuned extended LLM to 2048k via LongRoPE search}. In the final stage, we carry out an additional search procedure on the LLM that was previously fine-tuned at a context length of 256k. Through this process, the context window is substantially increased to 2048k, eliminating the need for extra fine-tuning. This represents an overall extension factor of $512\times$.

\textbf{Degradation in Original Context Window Performance}.  Upon expanding the context window to a substantial 2048k length, there is an observable decline in model performance inside the initial context window. This phenomenon, which is well-documented in positional interpolation research~\cite{pi}, arises because expanding the range compresses the positional embeddings within the standard context window into a more restricted space. This compression adversely influences the efficacy of the language model. Notably, when a 512$\times$ extension is applied, the embeddings corresponding to the original 4k context window are especially densely packed.

To address this issue, we conduct an additional evolution search on the extended LLM to fine-tune the RoPE rescale factors specifically for shorter context lengths, such as 4k and 8k. For these shorter contexts, the upper bound for the searched $\lambda$ is lowered, since there is a reduced need for positional interpolation. At inference time, the LLM automatically adapts the relevant RoPE rescale factors as appropriate.

\section{LongRoPE}

\label{sec:method}
Building upon these observations, we introduce LongRoPE. This approach leverages a novel, efficient search algorithm designed to maximize the advantages of the two identified non-uniformities. Utilizing this method, we extend the LLM context window to accommodate more than 2 million tokens.

\subsection{Problem Formulation}

These two forms of non-uniformity significantly expand the possible solution space, complicating the optimization process. To manage this, we recast the optimization of multidimensional non-uniform position interpolation into a search problem.

For a LLM targeting a  context window size of $L'$ and lengthy input documents $\mathbf{X}$, where each $\mathbf{x}\in\mathbf{X}$ surpasses $L'$ in token length, we denote the original rotary angle of the $i^{th}$ dimension in RoPE embedding at token position $n$ as  $\frac{n}{\beta^i}$. The optimization problem is then formulated as follows:

\begin{equation}
\begin{gathered}
		\label{eq:problem}

\underset {\mathbf{x}\in\mathbf{X}; \, |\mathbf{x}|\ge L'}{ \operatorname {arg\,min} } \,  \mathcal{L}\, \left( \text{LLM} (\text{RoPE}, \mathbf{X})\right), \quad \text{with}
\\
\scalebox{0.93}{
$\underset {\begin{subarray}{l} i=0,\cdots,\frac{d}{2}-1;\\ n\in[0,|\mathbf{x}|);\end{subarray}}{\text{RoPE}_(n)} = {\left[\cdots, \cos\left(\mathbb{I}(\hat{\lambda}_{i}, \hat{n})\frac{n}{\beta^i}\right), \sin\left(\mathbb{I}(\hat{\lambda}_{i}, \hat{n})\frac{n}{\beta^i}\right), \cdots\right] }$}
\\
\text{with} \quad \mathbb{I}(\hat{\lambda}_{i},\hat{n})=\begin{cases}
1 & \text{if } n < \hat{n} \\
\frac{1}{\lambda_{i}} & \text{if } n \geq \hat{n}
\end{cases}
\end{gathered}
\end{equation}

In this context, we define a collection of rescaling factors, $\mathbb{I}(\hat{\lambda}_{i},\hat{n})$, to address both types of non-uniformity present. The parameters $\hat{\lambda_i}$ and $\hat{n}$ represent the irregularities in RoPE dimensional scaling and the positions of tokens, respectively. More precisely, $\mathbb{I}(\hat{\lambda}_{i},\hat{n})$ is employed to adjust the rotation angle associated with the $i^{th}$ dimension in RoPE, where $\hat\lambda_{i}$ acts as the corresponding scaling parameter, and $\hat{n}$ defines the positional threshold for token indices. For the first $\hat{n}-1$ token positions, the scaling factor $\hat\lambda_{i}$ is omitted and the standard RoPE rotation angle $\frac{n}{\beta^i}$ remains unchanged. However, once the token position satisfies $n\ge\hat{n}$, the introduced scaling factor becomes effective.

Suppose we are given a desired context window of length $L'$. Our aim is to determine the set of optimal rescaling factors ($\mathbb{I}(\hat{\lambda}_{0},\hat{n})$, $\mathbb{I}(\hat{\lambda}_{1},\hat{n})$, ..., $\mathbb{I}(\hat{\lambda}_{i},\hat{n})$, ...) for each RoPE dimension, from the $1^{st}$ to the $d^{th}$. By appropriately rescaling the $\text{RoPE}$ in this manner, the target $\text{LLM}$ can minimize the next token prediction loss, $\mathcal{L}$ (i.e., perplexity), over input sequences $\mathbf{X}$ with token length $L'$.

\subsection{Searching the Non-uniform Position Interpolation}

To address the issue described in Eq.~\ref{eq:problem}, we propose an efficient and straightforward approach that identifies the optimal RoPE rescale factors, thereby taking full advantage of the complex, multidimensional variations present in position embedding.

\textbf{Search space}. To capture both forms of non-uniformity, we construct an expansive search space, as summarized in Table~\ref{tbl:searchspace}. Our approach permits the optimization of individual rescale factors for every dimension within the RoPE embedding. For ease of configuring the search space, we opt to search over $\lambda_i$ and $\hat{n}$, rather than directly searching $\mathbb{I}(\hat{\lambda}_{i},\hat{n})$, with the relationship $\hat{\lambda}_{i}=1/\lambda_i$. According to Table~\ref{tbl:searchspace}, the variable $\lambda_i$ spans a range starting from 1.0 (corresponding to straightforward extrapolation) and extends up to $s\times1.25$ (which allows for interpolation beyond that of PI), incremented in steps of 0.01. Here, $s$ represents the desired ratio for extending the context window.

The parameter $\hat{n}$ determines how many of the initial token positions preserve their original embeddings through RoPE, without applying position interpolation. In our experiments, $\hat{n}$ is chosen from the set \{0, 1, 2, 4, 8, 12, 16, 20, 24, 28, 32, 64, 128, 256\}. Setting $\hat{n}=0$ implies that every token position is modified using the optimized rescale factors.

\textbf{Evolution-based search}. The search space outlined in Table~\ref{tbl:searchspace} includes a vast array of positional interpolation options, which makes thorough exploration computationally infeasible. As an example, an extension factor of $s=4\times$ results in $400^{128/2}\times14=4\times10^{167}$ possible configurations. The size of this space increases exponentially as the extension ratio becomes larger. To efficiently navigate this enormous search space, we employ evolution search~\cite{guo2020single} and incorporate two optimization methods that considerably enhance the efficiency of the process. The general structure of our search strategy is presented in Algorithm~\ref{alg:search}.

\textit{Optimized initial population generation}. Rather than selecting all $P$ rescale factors at random to form the population, we explicitly include the three RoPE rescale factors representing PI, NTK, and YaRN as initial individuals. The other $P$-3 members are generated by applying random mutations with probability $p$ to these three baseline rescale factors.

\textit{Monotonically non-decreasing constraint}. Following the creation of the initial population, we evaluate the LLM perplexity for every individual. This involves applying each individual's associated RoPE rescale factors to the target LLM and then calculating the perplexity on input $\mathbf{X}$. The individuals with the highest performance—specifically, the top-k—are selected as parents for the subsequent evolutionary process. Nonetheless, due to the enormous size of the search space, standard mutation and crossover can often yield suboptimal candidates, resulting in excessive and redundant perplexity evaluations. This inefficiency is exacerbated when $L'$ is large, as each perplexity inference is computationally demanding.

To tackle this issue, we enforce a monotonic non-decreasing condition on the sampled RoPE rescaling factors: $\lambda_i\le \lambda_{i+1}$. Only those RoPE rescalings that conform to this requirement are considered for application in the LLM perplexity evaluations, which markedly decreases the cost of searching over possible parameter values. Concretely, we stipulate that the sequence of $\lambda_i$ values must be monotonically increasing relative to the RoPE dimension (for $i=0,\dots,63$). This requirement of dimensional monotonicity is supported by NTK theory~\cite{jacot2018neural,tancik2020fourier,ntk}, which holds that lower index dimensions—associated with higher frequencies—should have lower levels of interpolation (i.e., smaller $\lambda_i$ values), while higher index dimensions—with lower frequencies—are suited to greater interpolation (i.e., larger $\lambda_i$ values).

\begin{algorithm}[t]
	\small
	\caption{The search algorithm}
	\textbf{Input:} target LLM, input samples $\mathbf{X}$,   population size $P$, mutation size $N_1$, crossover size $N_2$, max iterations $\mathcal{T}$, mutate probability $p$\\

	\begin{algorithmic}[1]
		\label{alg:search}
		\STATE Initialize Top-k to $\phi$;	
		\STATE $\text{P}_0$ $\leftarrow$ \textit{Initialize\_population\_with\_optimization} ($P$, $\mathbf{X}$, $p$);
		\FOR{$i$ from $1$ to $\mathcal{T}$}
		\STATE \textit{Compute\_perplexity} (LLM, $\text{P}_{i-1}$, $\mathbf{X}$);
		\STATE  Top-k $\leftarrow$ \textit{Update\_Topk} (Top-k, $\text{P}_{i-1}$);
		\STATE $\text{P}_{mutation}$ $\leftarrow$ \textit{Mutation\_with\_mono\_constraint} (Top-k, $N_1$, $p$);
		\STATE $\text{P}_{crossover}$ $\leftarrow$ \textit{Crossover\_with\_mono\_constraint} (Top-k, $N_2$);
		\STATE $\text{P}_i$ $\leftarrow$ $\text{P}_{mutation} \cup \text{P}_{crossover} \cup$ Top-k;
		\ENDFOR
		\STATE Output the member of Top-k with the minimum perplexity;
	\end{algorithmic}
\end{algorithm}

\textbf{8$\times$ extension without fine-tuning}. Leveraging our evolutionary search strategy, we are able to pinpoint non-uniform RoPE rescale parameters that maintain crucial dimensions and token positions, thereby reducing information loss caused by interpolation. As shown in Fig.\ref{fig:non-ft}, this approach enables the context window of LLaMA2 to increase from 4k to 32k with no fine-tuning required. By comparison, alternative techniques like PI, along with non-uniform NTK and YaRN, exhibit a sharp increase in perplexity when the context length is merely doubled.

\subsection{Extending LLM Context Window to 2048K}
\label{sec:extendto2048k}

\textbf{Progressive extension to 2048k}. Here, we present our approach for expanding the context length of pre-trained LLMs beyond the standard 4k up to more than 2048k. Our results indicate that employing non-uniform positional interpolation allows for an 8$\times$ increase in context size without the need for fine-tuning. When aiming for substantially larger expansions (such as 512$\times$), however, fine-tuning becomes indispensable. One strategy involves identifying appropriate RoPE rescaling factors suitable for the extended 2048k window, followed by fine-tuning. Despite this, such a process is hindered by the substantial computational resources it requires. Additionally, our observations suggest that effectively fine-tuning LLMs at such high extension ratios remains challenging (refer to Appendix).

Fortunately, LongRoPE demonstrates efficacy when applied to both base and further fine-tuned extended LLMs. To this end, we propose a resource-efficient, incremental strategy that reaches the desired 2048k context length through merely 1k fine-tuning steps, all conducted at a training length capped at 256k.

$\Diamond$ \textit{LongRoPE search for scaling pre-trained LLMs to 256k contexts}. Using LLaMA2 as a reference, we perform a search to achieve context window lengths of 128k and 256k. At this phase, the enlargement factors correspond to 32$\times$ and 64$\times$, respectively.

$\Diamond$ \textit{Fine-tuning to 256k}. Next, we adapt the pre-trained LLM to support a 256k context window via a staged fine-tuning approach. We begin by fine-tuning LLaMA2 for 400 steps utilizing RoPE rescaled factors configured for 128k. Upon completion, we update the RoPE rescaled factors to 256k on the resulting checkpoint and proceed with an extra 600 fine-tuning steps. This two-stage process demonstrates greater efficiency compared to immediately fine-tuning for 256k.

$\Diamond$ \textit{Scaling a fine-tuned extended LLM to 2048k using LongRoPE search}. In the last step, we carry out an additional search process on the LLM that was previously fine-tuned for sequences of length 256k. This procedure successfully achieves a context window as large as 2048k, eliminating the need for any additional fine-tuning stages. The overall extension factor reaches $512\times$.

\textbf{Reduced performance in original context window}. Upon expanding the context window to an extensive 2048k length, a decrease in performance is observed within the initial context span. This outcome is a documented limitation of positional interpolation~\cite{pi}, which compresses position embeddings from higher-index tokens into a significantly more confined segment inside the original context window, thereby impairing the language model's effectiveness. Specifically, when the extension ratio reaches 512$\times$, position encodings within the prior 4k context window become notably congested.

To address this issue, an additional evolution search is conducted on the extended LLM, specifically to fine-tune the RoPE rescale factors at shorter context lengths such as 4k and 8k. Because these shorter lengths involve less positional interpolation, the search space for the maximum permitted $\lambda$ is constrained accordingly. At inference time, the LLM adapts the relevant RoPE rescale factors on-the-fly.

 \section{Experiments}

\label{sec:exp}
\subsection{Setup}
\textbf{Evaluation Tasks and Models}. Our experiments incorporate LongRoPE into LLaMA2-7B and Mistral-7B, with evaluations conducted across three primary areas: (1) perplexity measurements for extended-context LLMs handling lengthy documents; (2) a Passkey retrieval challenge, designed to assess the capability of a model to extract a simple passkey from a large body of distractor text; and (3) established LLM benchmark tests, constrained to a 4096-token context window.

\textbf{Fine-tuning}. In our approach with LLaMA2, we adopt a linearly decaying learning rate starting at 2e-5 and set the overall batch size to 32. The model undergoes fine-tuning for 400 steps on the Redpajama~\cite{redpajama} dataset, which is partitioned into 128k-length segments, each framed by BOS and EOS tokens. Following this, we continue training from the resulting checkpoint for an additional 600 steps to extend the context window to 256k. Training at the 128k context length is performed using 8 A100 GPUs within a distributed environment~\cite{cube}, whereas scaling up to 256k requires 16 A100 GPUs. For the Mistral model, we opt for a fixed learning rate of 1e-6 and utilize a global batch size of 64. Both the 128k and 256k context models adhere to the configuration described in YaRN~\cite{yarn}, consisting of 400 training steps on the Together Computer's Long-Data Collections~\cite{mistraldata} with a sequence length of 16k, using 4 A100 GPUs for all experiments.

\textbf{Search}. When the target window size does not exceed 256k, we set the parameters as follows: $P$=64, $N_1$=$N_2$=16, $p$=0.3, and $\mathcal{T}$=40, selecting the top 32 candidates for mutation and crossover at each iteration. For perplexity measurement, 5 random samples are drawn from the PG19 validation dataset, each satisfying a minimum length corresponding to the target context size. For target windows above 512k, the number of individuals in the population, as well as the counts for mutation and crossover, are reduced by half. In these cases, perplexity is evaluated on 3 randomly selected samples from the Pile-Books3~\cite{pile} validation set.

\textbf{Baselines}. To achieve a 2048k context window, we applied fine-tuning to models with initial 128k and 256k context lengths. These configurations are denoted as LongRoPE-2048k (ft=128k) for LLaMA2 and LongRoPE-2048k (ft=256k) for Mistral. We assess these four models against leading baseline approaches for extending context windows. In particular, we focus on open-source LLMs that underwent fine-tuning following positional interpolation using the PI, NTK, and YaRN methods. The baseline models considered in comparison include Together-32k~\cite{together}, Code LLaMA~\cite{codellama}, LongLoRA-full-FT-100k~\cite{longlora}, as well as YaRN-LLaMA and YaRN-Mistral~\cite{yarn}.

\subsection{Main Results}

\label{sec:mainresult}
\textbf{Language modeling on extended sequences up to 256k tokens}. To start, we assess our method against leading long-context LLMs for sequences of up to 256k tokens. Our analysis spans two benchmark datasets—Proof-pile~\cite{pg19} and the PG19~\cite{pile} test portions—to illustrate the breadth of our approach. We report perplexity across a range of context lengths, employing a sliding window with a size of 256 tokens. In the case of PG19, evaluation is performed using the entire test set, which comprises 100 documents. For Proof-pile, in alignment with the YaRN protocol~\cite{yarn}, we sample 10 documents at random, ensuring each selected sample is at least 128k tokens in length.

Table~\ref{tbl:proof-pile} and Table~\ref{tbl:pg19} present a perplexity comparison between LLaMA2 and Mistral, each expanded using several interpolation strategies, on Proof-pile and PG19 datasets, respectively. Two principal findings emerge: \textbf{(1)} the models extended with our approach consistently display reduced perplexity as the evaluation length increases from 4k to 256k, indicating enhanced utilization of extended context spans. \textbf{(2)} Notably, despite operating with a context window 16$\times$ greater than typical settings—a scenario that generally risks degradation of short-context performance—our LongRoPE-2048k variants surpass the latest competitive baselines for context lengths up to 256k.

\textbf{Long sequence language modeling beyond 2000k}. To assess performance on exceptionally lengthy documents, the Books3~\cite{pile} dataset is utilized. For the sake of efficient evaluation, 20 books are chosen at random, each with lengths greater than 2048k. A sliding window of 256k is applied for analysis.

Table~\ref{tbl:books3} demonstrates that LongRoPE effectively expands the context window of both LLaMA2-7B and Mistral-7B up to 2048k, while maintaining perplexity on par with or better than existing baselines in shorter ranges (8k–128k). Distinct disparities in performance emerge between LLaMA2 and Mistral at the 2048k length. While Mistral achieves superior results over the baselines for shorter contexts, its perplexity rises above 7 beyond 256k. In contrast, LLaMA2 displays anticipated behavior: perplexity gradually declines as the context length increases, though it rises slightly at 1024k and 2048k. Notably, for LLaMA2, using LongRoPE-2048k with a fine-tuning context of 256k leads to improved outcomes compared to fine-tuning at 128k; this is attributable to the reduced secondary extension ratio (i.e., 8$\times$ compared to 16$\times$). Conversely, Mistral benefits more from a 128k fine-tuning window. This is largely because, for Mistral's 128k and 256k fine-tuning, the YaRN approach is used with a training sequence length of 16k, constraining Mistral's capacity to further extend its context window following fine-tuning.

\textbf{Passkey retrieval}. In this section, we examine the practical context window size within generative tasks. Our approach replicates the synthetic passkey retrieval evaluation introduced by~\cite{passkey}. Specifically, this task requires the model to locate a randomly generated passkey (a five-digit code) embedded somewhere within an extended document. The prompt structure can be found in the appendix. For evaluation, we conduct the passkey retrieval task 10 times, each instance positioning the passkey randomly at a location uniformly sampled along the total evaluation context length.

As illustrated in Fig.~\ref{fig:passkey}, our model’s retrieval accuracy is compared against established baselines. The performance of previous models falls to zero once the sequence length exceeds 128k tokens. In stark contrast, LongRoPE-LLaMA2-2048k (ft=256k) consistently achieves high retrieval accuracy (at least 90\%) across a broad range from 4k to 2048k tokens, even under the demanding condition of extracting a passkey from a token set of millions. Similarly, LongRoPE-Mistral-2048k (ft=128k) sustains perfect (100\%) retrieval rates up to 1800k tokens before decreasing to 60\% at 2048k, a trend that corresponds with the modest rise in perplexity at 2048k observed in Table~\ref{tbl:books3}.

\textbf{Evaluation on standard benchmarks within the original context window}. We assess the performance of LongRoPE-2048k models on their initial context length, utilizing the Hugging Face Open LLM Leaderboard~\cite{huggingfaceboard} under both zero-shot and few-shot conditions. Our experiments adopt 25-shot ARC-Challenge~\cite{allenai:arc}, 10-shot HellaSwag~\cite{zellers2019hellaswag}, 5-shot MMLU~\cite{mmlu}, as well as 0-shot TruthfulQA~\cite{lin2021truthfulqa}.

Table~\ref{tbl:commonsense} indicates that our models perform on par with the original benchmark, which was constructed for models with shorter context windows, and surpass the original Mistral on TruthfulQA by a margin of +0.5\%. Although LongRoPE-LLaMA2-2048k, which was fine-tuned at 256k, demonstrates a modest decline in performance, it still delivers results within acceptable limits for the majority of evaluated tasks.

\subsection{Ablation Results}

\textbf{Effectiveness of the second positional interpolation}. As part of our progressive extension strategy, we employ our search algorithm to perform a second round of non-uniform positional interpolation on LLMs that have already undergone fine-tuning and extension. To assess the benefits of this approach, we conduct empirical evaluations on our fine-tuned LLaMA2-256k model, further scaling its context length to 512k, 1024k, and 2048k using both PI and YaRN techniques. According to the results presented in Table~\ref{tbl:secondsearch}, our approach with non-uniform positional interpolation maintains stable perplexity across the increased sequence lengths. In comparison, PI and YaRN exhibit sharp increases in perplexity as the extension ratio grows.

\textbf{Effectiveness of recovery at shorter context lengths.} In order to address the drop in performance observed at reduced context lengths, we modify the RoPE factors for LongRoPE-2048k by employing our search algorithm. In this process, we specifically lower the upper bound on scale factors during the search to reduce interpolation for sequences of 4k and 8k tokens. 
Table~\ref{tbl:shortperformance} presents perplexity measurements for LongRoPE-LLaMA2-2048k on Proof-pile at both 4k and 8k lengths, in addition to mean accuracy across various LLM benchmarks. These findings indicate that the proposed method leads to marked enhancements in performance when processing shorter context lengths.

\textbf{Analysis of the two types of non-uniformities}. To determine the specific impact of each non-uniformity, we conduct ablation studies. Two separate experiments are designed: (i) we expand LLaMA2-7B to shorter context lengths of 16k and 32k with multiple strategies—including PI, searching exclusively over the RoPE dimension, and jointly optimizing both non-uniformity aspects; (ii) using the same approach, we take the previously fine-tuned 256k-length LLaMA2 model and further extend it to 2048k. Evaluation of perplexity is performed directly, without any subsequent fine-tuning. As reported in Table~\ref{tbl:searchdimension}, leveraging non-uniformity within the RoPE dimension leads to a marked perplexity reduction compared to using linear interpolation in PI. Introducing non-uniformity in token position yields performance gains at 16k and 32k, whereas no clear benefit is observed for the much longer 2048k context length, possibly reflecting the limits of this method at extreme scales. Simply retaining the earliest tokens without any interpolation does not appear helpful, suggesting an area for future exploration.

\section{Related Works}

Apart from techniques relying on position interpolation, this section examines other relevant approaches in the literature.

\textbf{Retrieval-based} methods incorporate an external memory system that stores long-term contextual information, leveraging retrieval mechanisms to access pertinent documents during inference~\cite{longllama,wang2023augmenting,borgeaud2022improving}. Implementing such methods generally requires substantial architectural changes to LLMs. In contrast, our approach introduces minimal alterations by slightly adjusting the positional embeddings. Furthermore, our method extends to a broader set of long-context scenarios beyond retrieval alone, including applications like comprehensive document summarization and few-shot learning.

\textbf{Attention-based context window extensions}. In addition to approaches like positional embedding interpolation, several studies extend the input context window of original LLMs by altering attention mechanisms within the existing context length~\cite{han2023lminfinite, streamingllm,parallelwindow}. Central to these methods is the use of innovative attention masks, which help address the problem of attention explosion from novel token positions. Notably, such techniques are complementary to strategies involving positional interpolation.

\textbf{Fine-tuning based approaches} investigate strategies to efficiently adapt pre-trained LLMs with adjusted position embeddings to support extended context lengths. Prior research, such as Code LLaMA~\cite{codellama}, LLaMA2 Long~\cite{xiong2023effective}, and ScaledRoPE~\cite{liu2023scaling}, adopts extremely large base values for RoPE and conducts fine-tuning at the desired sequence lengths. Our approach, in contrast, is versatile across a range of target lengths, supporting sequences exceeding 2M tokens. As fine-tuning for extremely long contexts (e.g., above 128k) poses significant GPU demands, methods like LongLoRA~\cite{longlora} and PoSE~\cite{zhu2023pose} have emerged to alleviate these computational challenges. Our technique remains complementary and independent of these resource-efficient fine-tuning strategies.

\section{Conclusion}

In this paper, we introduce LongRoPE, a technique that dramatically increases the context window of LLMs to an unprecedented 2048k tokens, all while preserving their performance within the original, smaller context lengths. Our approach leverages two distinct types of non-uniformities found in RoPE positional encodings, identified through an efficient evolutionary search process. This strategy delivers two main advantages: it supplies advantageous initialization for subsequent fine-tuning and permits an 8$\times$ expansion of the context window even without additional fine-tuning. Building upon this foundation, we outline a gradual extension approach that utilizes LLMs fine-tuned on 256k-length contexts to achieve a 2048k context length, without necessitating further fine-tuning. Comprehensive experiments demonstrate the efficacy of LongRoPE. We anticipate that our LongRoPE-2048k models will facilitate a broad range of long-context applications and spark new research directions in this area.

\section*{Broader Impacts} The research detailed in this paper aims to contribute to progress within the Machine Learning community. While the developments may have a range of societal implications, we do not identify any particular impacts that warrant emphasis at this time.

	\balance

\end{document}